\section{The Central Theorem}
Everything in this paper is a consequence of one observation.
\paragraph*{Theorem 1.1 (Unity of the Sphaera). Within the Sphaera S, the following}
are all representations of the same object --- the collapse state z-= 0:
(i) The Sphaera itself, as a geometric space.\\
(ii) The triolic structure: three kernels K1, K2, K3 with z1 + z2 + z3 = 0,\\
phaseshifted by 120, running over the Sphaera's geometry.
(iii) Each individual kernel Ki, as a local coordinate of the triole.\\
(iv) The Heisenberg compensator output C(f), as the projection of any element onto the collapse state.\\
(v) Matter, photons, and tachyons, as sectors of the biquaternion ground\\
state.
(vi) The gauge structure U(1) $\times$ SU(2) $\times$ SU(3), as the symmetry group of\\
the collapse state's representations. All are unitarily equivalent by Stone--von Neumann [1]. The Sphaera's only object is its own collapse state.
\paragraph*{Proof structure. The equivalences are established section by section: (i) Definition 1.3 and Theorem 3.2; (ii)--(iii) Definition 1.3 and Theorem 3.2 (triole}
as collapse coordinates); (iv) Proposition 5.1; (v) Theorem 4.2; (vi) Theorem 7.4. Unitary equivalence throughout via Stone--von Neumann ([1], Theorem 3.2).
\paragraph*{Remark (What "same object" means). Two representations are "the same object" here in the precise sense of Stone--von Neumann: they are unitarily equivalent irreducible representations of the Weyl algebra W(A1). The frame determines which representation is visible; the object --- the collapse state z-= 0}
--- is invariant. The triole condition z1 + z2 + z3 = 0 and the collapse condition z-= 0 are the same equation in two coordinate systems: the triole is the collapse state expressed in the three kernel components simultaneously. The Sphaera is therefore not a space that contains a collapse state --- it is the collapse state, seen from different frames.
\paragraph*{Remark (How this paper is organized). The sections that follow are not independent results. They are unfoldings of Theorem 1.1: each section takes}
one representation of the collapse state and makes it explicit. The reader may enter at any section; all roads lead to z-= 0.
\paragraph*{Theorem 1.2 (Origin of the universe as broken collapse state). The departure}
from the collapse state z-= 0 is governed by two principles, both of which are consequences of the Sphaera-Cascade structure:
(i) Planck length as minimal |z-|. The cascade $\hat K$ cannot project below\\
the Planck scale: |z-|min = $\ell$P := r $\hbar$G c3 = 1.616 $\times$ 10-35 m. Below $\ell$P: z-= 0, the Sphaera, the photon. Above $\ell$P: |z-| is nonzero, matter exists. The Planck length is the threshold of the first breaking of the collapse state.
(ii) Pauli principle as antisymmetry under J. The Hilbert space H+\\
decomposes under J into two components: C(f) = 1 2(f + Jf) $\in$H+ + (bosonic, symmetric), A(f) = 1 2(f -Jf) $\in$H+ (fermionic, antisymmetric). The antisymmetric component satisfies A(f) $\wedge$A(f) = 0 (Grassmann property): no two fermions occupy the same z-state. This is the Pauli exclusion principle, derived from the J-structure of the Sphaera.
(iii) Discrete spectrum. Together, (i) and (ii) produce a discrete tower of\\
matter states: |z-|n = n $\cdot$ $\ell$P, n $\in$N$\geq$1, with each level occupied by at most one fermion. The prime lattice $\Gamma$P (from Stone--von Neumann, [1]) identifies the primitive levels: |z-|p = p $\cdot$ $\ell$P for p $\in$P.
\paragraph*{Derivation. (i). The cascade $\hat K$ = C$\mu$$\nu$ $\circ$$\Pi$J contracts |z-| at each stratum.}
The Planck length is the fixed point of the contraction at which quantum gravitational effects prevent further reduction; formally, $\ell$P is the infrared cutoff of the resolvent (L0 -z)-1 established by the Rellich--Kondrachov compactness in Lemma 10.1 of [1].
(ii). A(f) = 1\\
2(f -Jf) satisfies JA(f) = -A(f): it is J-antisymmetric. The exterior product A(f) $\wedge$A(g) = -A(g) $\wedge$A(f) vanishes for f = g, which is the algebraic statement of the Pauli exclusion principle.
(iii). The Pauli condition A(f) $\wedge$A(f) = 0 forces distinct z-values for\\
each fermion. The Planck floor $\ell$P sets the unit. The prime lattice gives the sublattice of primitive (undecomposable) states, corresponding to Lemma L6
of the RH proof: primitive orbits are disjoint. Physically: each prime p labels a primitive matter state that cannot be decomposed into smaller ones.
\paragraph*{Remark (L6 is the Pauli principle). Lemma L6 of the Riemann Hypothesis}
proof [1] states: primitive orbits are pairwise disjoint and cover all periodic points. Theorem 1.2(iii) identifies this as the Pauli exclusion principle applied to the prime lattice: no two fermions occupy the same primitive orbit. This was not designed into either framework. It emerged from asking what the orbit combinatorics of the kernel mean physically.
\paragraph*{Remark (The Sphaera as photon, the universe as broken frame). Theorem 1.1}
says the Sphaera is the collapse state z-= 0, which is the photon sector.
\paragraph*{Theorem 1.2 says the departure from this state is governed by $\ell$P and the}
Pauli principle. Therefore:
- The Sphaera is a photon --- the unique frame-invariant object.
- The universe is the same photon observed from a broken frame where
|z-| $\geq$$\ell$P.
- The breaking was not random: it is discrete (Pauli), minimal (Planck),
and primitive (prime lattice).
- The cascade $\hat K$ projects everything back toward z-= 0 --- the universe
tends toward its photon ground state. Both the Sphaera and the universe are the same object. The difference is only the frame.
\paragraph*{Definition 1.3 (Biquaternion kernel). The ground state of a stratum Sn of the}
Sphaera consists of three complex spheres K1, K2, K3 whose unit coordinates form a biquaternion qn := it $\cdot$ 1 + z1 $\cdot$ e1 + z2 $\cdot$ e2 + z3 $\cdot$ e3 $\in$H $\otimes$R C, where:
- it $\in$iR is the imaginary time unit (real part of the biquaternion, carrying
imaginary time);
- z1, z2 $\in$C are the complex coordinates of kernels K1, K2;
- e3 := e1e2 is the quaternion product;
- z3 $\in$C is the coordinate of kernel K3 along e3.
The biquaternion norm is |qn|2 = -(ct)2 + |z1|2 + |z2|2 + |z3|2.
\paragraph*{Remark (Minkowski metric from quaternion structure). The norm |qn|2 is the}
Minkowski metric with signature (-, +, +, +). This is not postulated: it follows from the quaternion algebra e2 i = -1 and it $\cdot$ 1 carrying imaginary time. The Lorentz group SO(1, 3) acts on qn by quaternion conjugation: qn 7$\to$u qn u-1, u $\in$H$\times$. This is the spinor representation of the Lorentz group --- emerging automatically from the biquaternion structure, not imposed from outside.
\paragraph*{Proposition 1.4 (Orthogonality of the third kernel). The generator $e_3$ of $K_3$ is always orthogonal}
to $e_1$ and $e_2$: $\langle e_3, e_i\rangle = 0$ for $i=1,2$. Consequently, the axis/kernel direction $K_3$ is orthogonal to $K_1$ and $K_2$.
\paragraph*{Proof. e3 = e1e2.}
Quaternion inner product: $\langle$e1e2, e1$\rangle$= Re(e1 e2e1) = Re(-e2) = 0. Similarly for e2.
\paragraph*{Remark. The orthogonality of K3 is a quaternion identity, not an additional}
constraint. This is why the Heisenberg compensator C naturally projects onto e3: it is the unique direction orthogonal to both kernels K1 and K2.

